\section{Conclusions}\label{sec:conclusion}

A hidden attractor can be found faster than by uniform sampling only through information that the trajectories missing the target carry about it, and \cref{thm:info} quantifies this: the guaranteed gain is at most about $2.6/\mu(Q)$, where $Q$ is the set around the target on which the observed information is constant. For rare hidden regimes of systems of $k$ units, escape times make this set exponentially small in $k$, and ideal multilevel splitting along the escape time attains the bound up to a constant (\cref{prop:product}), so the bound cannot be improved by more than this constant. The portfolio of splitting along transient duration, active learning of basin boundaries, and a uniform share realizes this potential within a factor of 1.5--2 per unit on coupled copies, where its gain grows from 65 to $1.3\times10^4$ for two to four copies; it changes the cost from $p^{-1}$ to $p^{-0.45}$ where the transient signal is present on all scales, is 3--7 times faster on the rare hidden attractors of the benchmark that have such a signal, and is guaranteed never to be more than $1/\theta$ times costlier than uniform sampling in integrations, with every autonomous arm, such as stagger-and-step chains, contributing its order. Where trajectories are bounded and escape times are worthless, the settling part of the transient duration still gives gains of 7--360, and they carry over to wall time.

The benchmark, the identification and verification protocol, and the synthetic families with exactly known basin measure are independent of our method and can serve as a common test bed: the cost $N_{0.9}$ relative to uniform sampling is a single number that any localization method can report. Beyond the method, the work leaves the hiddenness measure $h_\varepsilon$, which grades hidden attractors by the scale at which the equilibria cease to lead to them, and a hidden chaotic attractor of two coupled Chua circuits that a single circuit does not have.
